\section{Formal setting (OC Core / ETS only)}

This section fixes the formal setting used throughout the paper.
All notions are taken exclusively from the Ontology of Continua (OC Core)
and its executable specification ETS.K\_EA.
No new primitives, semantic extensions, or external interpretations are
introduced.

\medskip

\begin{definition}[Enterprise continuum]
The \emph{enterprise continuum} is the formal object
\[
K_{EA} \;=\; (\Omega(K_{EA}), A, P, \Theta, J, C, k)
\]
as defined in OC Core and instantiated by ETS.K\_EA.
Here, \(\Omega(K_{EA})\) denotes the space of admissible enterprise states,
while \(A\), \(P\), \(\Theta\), \(J\), \(C\), and \(k\) denote axes,
potentials, thresholds, flows, cycles, and the continuum measure,
respectively.
All admissibility conditions, invariants, and structural constraints are
fixed by the executable specification.
\end{definition}

\medskip

\begin{definition}[Model realization]
A \emph{model realization} is any state
\[
s \in \Omega(K_{EA})
\]
that satisfies all axioms, invariants, threshold conditions, and structural
constraints imposed by OC Core and enforced by ETS.K\_EA.
No assumptions are made regarding uniqueness, identifiability, or
observability of such states.
\end{definition}

\medskip

\begin{definition}[Observation language]
The \emph{observation language} is the restricted formal vocabulary exposed
by ETS.K\_EA for expressing statements about enterprise states.
It includes only quantities, relations, predicates, and constraints that
are explicitly representable within the executable model.
Statements formulated outside this language are not admissible inputs to
the model.
\end{definition}

\medskip

\begin{definition}[Admissible observation]
An \emph{admissible observation} is a finite set of well-formed statements in
the observation language whose joint interpretation does not presuppose the
violation of any axiom, invariant, threshold, or structural constraint
defining \(\Omega(K_{EA})\).
Admissibility is a property determined prior to any diagnostic procedure or
falsifiability assessment.
\end{definition}

\medskip

\begin{definition}[Compatibility]
An admissible observation \(O\) is \emph{compatible} with a model realization
\(s \in \Omega(K_{EA})\) if all statements in \(O\) can be embedded into
\(s\) without violating any mandatory invariant, threshold, or structural
constraint.
Compatibility does not require uniqueness of the realization and may hold
for multiple distinct states.
\end{definition}

\medskip

\begin{definition}[STOP outcome]
A \emph{STOP} outcome occurs when no admissible continuation of diagnosis,
state refinement, or inference exists that preserves compatibility with at
least one model realization in \(\Omega(K_{EA})\).
STOP is defined as a terminal condition of the diagnostic procedure and does
not constitute a statement about the emptiness or inconsistency of
\(\Omega(K_{EA})\).
\end{definition}

\medskip

\begin{remark}
This formal setting explicitly separates three layers:
(i) the existence of admissible model realizations in
\(\Omega(K_{EA})\);
(ii) the expressibility and admissibility of observations formulated in the
observation language; and
(iii) the behavior and termination conditions of diagnostic procedures,
including STOP.
All subsequent results on falsifiability and limits operate strictly within
this interface and rely on no external semantics.
\end{remark}
